\documentclass[10pt,a4paper]{amsart}
\usepackage[margin=19mm]{geometry}
\usepackage{amsmath,amssymb,mathtools}
\usepackage[scr=rsfs,cal=euler]{mathalfa}
\usepackage{xfrac}
\usepackage[unicode,hidelinks,bookmarks=false]{hyperref}
\newcommand{\ie}{i.e.\ }
\newcommand{\cf}{cf.\ }
\newcommand{\resp}{respectively\ }
\newcommand{\Subsection}{\subsection}
\providecommand{\nobreakdash}{\nobreak\hbox{-}}
\setlength{\emergencystretch}{2em}
\multlinegap0pt
\usepackage{bm}
\usepackage{amssymb}
\usepackage[normalem]{ulem}
\usepackage{xcolor}
\newtheorem{lemma}{Lemma}
\title{Watanabe--Strogatz Invariants in the Liouvillian Dynamics of Coupled Phase Oscillators via the Koopman Framework}
\author{Keisuke Taga, Hiroya Nakao}
\date{Source: arXiv:2603.18809v1 (CC BY 4.0)}
\begin{document}
\maketitle

\noindent\emph{Source and licence.} Watanabe--Strogatz Invariants in the Liouvillian Dynamics of Coupled Phase Oscillators via the Koopman Framework. Version arXiv:2603.18809v1 (CC BY 4.0), distributed by arXiv under the Creative Commons Attribution 4.0 International licence, http://creativecommons.org/licenses/by/4.0/, as recorded in the arXiv licence metadata for this exact version and shown on the abstract page https://arxiv.org/abs/2603.18809v1. Full licence notice: \texttt{licenses/arxiv-2603.18809-CC-BY-4.0.txt}.\par\smallskip

\noindent\textit{Editorial scope.} Counted unit: the article's lemma lem:decomposition (source0.tex lines 1143--1162). The article's complete original proof of it is delivered with it (source0.tex lines 1164--1280, one \texttt{proof} environment of 117 lines). No mathematics is written, repaired or re-derived: the statement and the proof are the article's own lines, in the article's own notation, and no retained line is substituted or re-flowed.  Retained uncounted as the source-local context the proof needs: lines 287--290.  Nothing was omitted from any retained span, so the package carries no omission record.  No retained line contains a citation, so the wrapper carries no bibliography and no bibliography entry.  No retained line contains a figure environment, an image-inclusion command or drawing code, so no image is delivered and the package declares no asset dependency.  Editorial formatting only, in addition: an AMS \texttt{amsart} preamble that restates the article's own declarations and macros used by the retained lines, namely the environments \texttt{lemma} on separate counters, together with the standard packages those lines need.  The retained environments print in this document as Lemma 1 in order of appearance, whereas the article prints the same items with the numbering of its own full document.  The complete native source of the article as distributed in the arXiv e-print for this exact version is bundled unchanged as \texttt{sources/196725/source0.tex}.
\begin{align}
    \label{eq:eigenfunction}
	\psi_{\bm{q}}^N(\bm{\theta}_N) := \frac{1}{\prod_{j=1}^N \sin\left( \frac{\theta_{q^{(j)}} - \theta_{q^{(j+1)}}}{2} \right)}.
\end{align}

\begin{lemma}
\label{lem:decomposition}
    Let $N \geq 5$, and let $\Psi^N_{\bm{q}_1, \bm{q}_2}$ be defined as before. Then $\Psi^N_{\bm{q}_1, \bm{q}_2}$ can be reduced to $\Psi^{N-1}_{\bm{q}'_1, \bm{q}'_2}$ or decomposed as a product of $\Psi^{N-1}_{\bm{q}'_1, \bm{q}'_2}$ and invariants $\langle  i,j,k,l\rangle$.
    Here,    
    $\bm{q}'_1$, $\bm{q}'_2$ are the permutations of $\{1,2,\ldots,N-1\}$ obtained by removing $N$ 
    from $\bm{q}_1$ and $\bm{q}_2$:
    \begin{align}
        q_1'^{(l)} &=
        \begin{cases}
            q_1^{(l)}, & \text{if } l < \xi, \\
            q_1^{(l+1)}, & \text{if } l \geq \xi,
        \end{cases} \quad
        q_2'^{(l)} =
        \begin{cases}
            q_2^{(l)}, & \text{if } l < \zeta, \\
            q_2^{(l+1)}, & \text{if } l \geq \zeta,
        \end{cases}
    \end{align}
    where $q_1^{(\xi)} = q_2^{(\zeta)} = N$, and we impose the cyclic property in the indices as $q_{1,2}^{(0)}=q_{1,2}^{(N)}$ and $q_{1,2}^{(N+1)}=q_{1,2}^{(1)}$.
\end{lemma}

\begin{proof}
%
From Eq.~\eqref{eq:eigenfunction}, we can decompose $\Psi^N_{\bm{q}_1, \bm{q}_2}$ as
%
\begin{align}
    \Psi^N_{\bm{q}_1, \bm{q}_2}
    = \Psi^{N-1}_{\bm{q}'_1, \bm{q}'_2}\frac{\sin\left( \frac{\theta_{q_1^{(\xi-1)}} - \theta_{q_1^{(\xi)}}}{2} \right)\sin\left( \frac{\theta_{q_1^{(\xi)}} - \theta_{q_1^{(\xi+1)}}}{2} \right)}
           {\sin\left( \frac{\theta_{q_2^{(\zeta-1)}} - \theta_{q_2^{(\zeta)}}}{2} \right)\sin\left( \frac{\theta_{q_2^{(\zeta)}} - \theta_{q_2^{(\zeta+1)}}}{2} \right)}\cdot \frac{\sin\left( \frac{\theta_{q_2^{(\zeta-1)}} - \theta_{q_2^{(\zeta+1)}}}{2} \right)}{\sin\left( \frac{\theta_{q_1^{(\xi-1)}} - \theta_{q_1^{(\xi+1)}}}{2} \right)}.
\end{align}
%
We analyze this by dividing into the following six cases. 
%
\begin{enumerate}
\item[\textit{Case.1:}] $\{{q_1^{(\xi-1)}},{q_1^{(\xi+1)}}\} = \{{q_2^{(\zeta-1)}}, {q_2^{(\zeta+1)}}\}$\\ \\
In this case, the following equation holds: 
\begin{align}
    \frac{\sin\left( \frac{\theta_{q_1^{(\xi-1)}} - \theta_{q_1^{(\xi)}}}{2} \right)\sin\left( \frac{\theta_{q_1^{(\xi)}} - \theta_{q_1^{(\xi+1)}}}{2} \right)}
           {\sin\left( \frac{\theta_{q_2^{(\zeta-1)}} - \theta_{q_2^{(\zeta)}}}{2} \right)\sin\left( \frac{\theta_{q_2^{(\zeta)}} - \theta_{q_2^{(\zeta+1)}}}{2} \right)} = 1.
\end{align}
Thus, we obtain
\begin{align}
    \Psi^N_{\bm{q}_1, \bm{q}_2} = \Psi^{N-1}_{\bm{q}'_1, \bm{q}'_2}.
\end{align}

This includes both the aligned case $(q_1^{(\xi-1)},q_1^{(\xi+1)})=(q_2^{(\zeta-1)},q_2^{(\zeta+1)})$ and the swapped case $(q_1^{(\xi-1)},q_1^{(\xi+1)})=(q_2^{(\zeta+1)},q_2^{(\zeta-1)})$, and in either situation the sine factors cancel pairwise, yielding the ratio $1$.

\item[\textit{Case.2:}] $q_1^{(\xi-1)} = q_2^{(\zeta-1)} = M$, ${q_1^{(\xi+1)}} \ne {q_2^{(\zeta+1)}}$ \\ \\
In this case, the following relation holds:
%
\begin{align}
    &\frac{\sin\left( \frac{\theta_{q_1^{(\xi-1)}} - \theta_{q_1^{(\xi)}}}{2} \right)\sin\left( \frac{\theta_{q_1^{(\xi)}} - \theta_{q_1^{(\xi+1)}}}{2} \right)\sin\left( \frac{\theta_{q_2^{(\zeta-1)}} - \theta_{q_2^{(\zeta+1)}}}{2} \right)}
   {\sin\left( \frac{\theta_{q_2^{(\zeta-1)}} - \theta_{q_2^{(\zeta)}}}{2} \right)\sin\left( \frac{\theta_{q_2^{(\zeta)}} - \theta_{q_2^{(\zeta+1)}}}{2} \right)\sin\left( \frac{\theta_{q_1^{(\xi-1)}} - \theta_{q_1^{(\xi+1)}}}{2} \right)}\cr
    &=\frac{\sin\left( \frac{\theta_{q_1^{(\xi)}} - \theta_{q_1^{(\xi+1)}}}{2} \right)\sin\left( \frac{\theta_{q_2^{(\zeta-1)}} - \theta_{q_2^{(\zeta+1)}}}{2} \right)}
   {\sin\left( \frac{\theta_{q_2^{(\zeta)}} - \theta_{q_2^{(\zeta+1)}}}{2} \right)\sin\left( \frac{\theta_{q_1^{(\xi-1)}} - \theta_{q_1^{(\xi+1)}}}{2} \right)}\cr
   &=\frac{\sin\left( \frac{\theta_{N} - \theta_{q_1^{(\xi+1)}}}{2} \right)\sin\left( \frac{\theta_{M} - \theta_{q_2^{(\zeta+1)}}}{2} \right)}
   {\sin\left( \frac{\theta_{N} - \theta_{q_2^{(\zeta+1)}}}{2} \right)\sin\left( \frac{\theta_{M} - \theta_{q_1^{(\xi+1)}}}{2} \right)}\cr
   &= \langle N,q_1^{(\xi+1)},M, q_2^{(\zeta+1)}\rangle.
\end{align}
Thus, we obtain
\begin{align}
    \Psi^N_{\bm{q}_1, \bm{q}_2} = \Psi^{N-1}_{\bm{q}'_1, \bm{q}'_2}\langle N,q_1^{(\xi+1)},M, q_2^{(\zeta+1)}\rangle.
\end{align}

\item[\textit{Case.3:}] ${q_1^{(\xi+1)}} = {q_2^{(\zeta+1)}} = M$, ${q_1^{(\xi-1)}} \ne {q_2^{(\zeta-1)}}$\\ \\
%
In this case, the following relation holds:
\begin{align}
    &\frac{\sin\left( \frac{\theta_{q_1^{(\xi-1)}} - \theta_{q_1^{(\xi)}}}{2} \right)\sin\left( \frac{\theta_{q_1^{(\xi)}} - \theta_{q_1^{(\xi+1)}}}{2} \right)\sin\left( \frac{\theta_{q_2^{(\zeta-1)}} - \theta_{q_2^{(\zeta+1)}}}{2} \right)}
   {\sin\left( \frac{\theta_{q_2^{(\zeta-1)}} - \theta_{q_2^{(\zeta)}}}{2} \right)\sin\left( \frac{\theta_{q_2^{(\zeta)}} - \theta_{q_2^{(\zeta+1)}}}{2} \right)\sin\left( \frac{\theta_{q_1^{(\xi-1)}} - \theta_{q_1^{(\xi+1)}}}{2} \right)}\cr
    &=\frac{\sin\left( \frac{\theta_{q_1^{(\xi-1)}} - \theta_{q_1^{(\xi)}}}{2} \right)\sin\left( \frac{\theta_{q_2^{(\zeta-1)}} - \theta_{q_2^{(\zeta+1)}}}{2} \right)}
   {\sin\left( \frac{\theta_{q_2^{(\zeta-1)}} - \theta_{q_2^{(\zeta)}}}{2} \right)\sin\left( \frac{\theta_{q_1^{(\xi-1)}} - \theta_{q_1^{(\xi+1)}}}{2} \right)}\cr
    &=\frac{\sin\left( \frac{\theta_{q_1^{(\xi-1)}} - \theta_{N}}{2} \right)\sin\left( \frac{\theta_{q_2^{(\zeta-1)}} - \theta_{M}}{2} \right)}
   {\sin\left( \frac{\theta_{q_2^{(\zeta-1)}} - \theta_{N}}{2} \right)\sin\left( \frac{\theta_{q_1^{(\xi-1)}} - \theta_{M}}{2} \right)}\cr
    &= \langle N,q_1^{(\xi-1)},M, q_2^{(\zeta-1)}\rangle.
\end{align}
Thus, we obtain
\begin{align}
    \Psi^N_{\bm{q}_1, \bm{q}_2} = \Psi^{N-1}_{\bm{q}'_1, \bm{q}'_2}\langle N,q_1^{(\xi-1)},M, q_2^{(\zeta-1)}\rangle.
\end{align}

\item[\textit{Case.4:}] ${q_1^{(\xi-1)}} = {q_2^{(\zeta+1)}} = M$, ${q_1^{(\xi+1)}} \ne {q_2^{(\zeta-1)}}$\\ \\
In this case, the following relation holds:
\begin{align}
    &\frac{\sin\left( \frac{\theta_{q_1^{(\xi-1)}} - \theta_{q_1^{(\xi)}}}{2} \right)\sin\left( \frac{\theta_{q_1^{(\xi)}} - \theta_{q_1^{(\xi+1)}}}{2} \right)\sin\left( \frac{\theta_{q_2^{(\zeta-1)}} - \theta_{q_2^{(\zeta+1)}}}{2} \right)}
   {\sin\left( \frac{\theta_{q_2^{(\zeta-1)}} - \theta_{q_2^{(\zeta)}}}{2} \right)\sin\left( \frac{\theta_{q_2^{(\zeta)}} - \theta_{q_2^{(\zeta+1)}}}{2} \right)\sin\left( \frac{\theta_{q_1^{(\xi-1)}} - \theta_{q_1^{(\xi+1)}}}{2} \right)}\cr
    &=-\frac{\sin\left( \frac{\theta_{q_1^{(\xi)}} - \theta_{q_1^{(\xi+1)}}}{2} \right)\sin\left( \frac{\theta_{q_2^{(\zeta-1)}} - \theta_{q_2^{(\zeta+1)}}}{2} \right)}
   {\sin\left( \frac{\theta_{q_2^{(\zeta-1)}} - \theta_{q_2^{(\zeta)}}}{2} \right)\sin\left( \frac{\theta_{q_1^{(\xi-1)}} - \theta_{q_1^{(\xi+1)}}}{2} \right)}\cr
    &=-\frac{\sin\left( \frac{\theta_{N} - \theta_{q_1^{(\xi+1)}}}{2} \right)\sin\left( \frac{\theta_{q_2^{(\zeta-1)}} - \theta_{M}}{2} \right)}
   {\sin\left( \frac{\theta_{q_2^{(\zeta-1)}} - \theta_{N}}{2} \right)\sin\left( \frac{\theta_{M} - \theta_{q_1^{(\xi+1)}}}{2} \right)}\cr
    &= -\langle N,q_1^{(\xi+1)},M, q_2^{(\zeta-1)}\rangle.
\end{align}
Thus, we obtain
\begin{align}
    \Psi^N_{\bm{q}_1, \bm{q}_2} = -\Psi^{N-1}_{\bm{q}'_1, \bm{q}'_2}\langle N,q_1^{(\xi+1)},M, q_2^{(\zeta-1)}\rangle.
\end{align}

\item[\textit{Case.5:}] ${q_1^{(\xi+1)}} = {q_2^{(\zeta-1)}} = M$, ${q_1^{(\xi-1)}} \ne {q_2^{(\zeta+1)}}$\\ \\
%
In this case, the following relation holds:
\begin{align}
    &\frac{\sin\left( \frac{\theta_{q_1^{(\xi-1)}} - \theta_{q_1^{(\xi)}}}{2} \right)\sin\left( \frac{\theta_{q_1^{(\xi)}} - \theta_{q_1^{(\xi+1)}}}{2} \right)\sin\left( \frac{\theta_{q_2^{(\zeta-1)}} - \theta_{q_2^{(\zeta+1)}}}{2} \right)}
   {\sin\left( \frac{\theta_{q_2^{(\zeta-1)}} - \theta_{q_2^{(\zeta)}}}{2} \right)\sin\left( \frac{\theta_{q_2^{(\zeta)}} - \theta_{q_2^{(\zeta+1)}}}{2} \right)\sin\left( \frac{\theta_{q_1^{(\xi-1)}} - \theta_{q_1^{(\xi+1)}}}{2} \right)}\cr
    &=-\frac{\sin\left( \frac{\theta_{q_1^{(\xi-1)}} - \theta_{q_1^{(\xi)}}}{2} \right)\sin\left( \frac{\theta_{q_2^{(\zeta-1)}} - \theta_{q_2^{(\zeta+1)}}}{2} \right)}
   {\sin\left( \frac{\theta_{q_2^{(\zeta)}} - \theta_{q_2^{(\zeta+1)}}}{2} \right)\sin\left( \frac{\theta_{q_1^{(\xi-1)}} - \theta_{q_1^{(\xi+1)}}}{2} \right)}\cr
    &=-\frac{\sin\left( \frac{\theta_{q_1^{(\xi-1)}} - \theta_{N}}{2} \right)\sin\left( \frac{\theta_{M} - \theta_{q_2^{(\zeta+1)}}}{2} \right)}
   {\sin\left( \frac{\theta_{N} - \theta_{q_2^{(\zeta+1)}}}{2} \right)\sin\left( \frac{\theta_{q_1^{(\xi-1)}} - \theta_{M}}{2} \right)}\cr
    &= -\langle N,q_1^{(\xi-1)},M, q_2^{(\zeta+1)}\rangle.
\end{align}
Thus, we obtain
\begin{align}
    \Psi^N_{\bm{q}_1, \bm{q}_2} = -\Psi^{N-1}_{\bm{q}'_1, \bm{q}'_2}\langle N,q_1^{(\xi-1)},M, q_2^{(\zeta+1)}\rangle.
\end{align}

\item[\textit{Case.6:}] $\{{q_1^{(\xi-1)}},{q_1^{(\xi+1)}}\} \cap \{{q_2^{(\zeta-1)}}, {q_2^{(\zeta+1)}}\} = \varnothing$ 
\\ \\
%
In this case, the following relation holds:
\begin{align}
    &\frac{\sin\left( \frac{\theta_{q_1^{(\xi-1)}} - \theta_{q_1^{(\xi)}}}{2} \right)\sin\left( \frac{\theta_{q_1^{(\xi)}} - \theta_{q_1^{(\xi+1)}}}{2} \right)\sin\left( \frac{\theta_{q_2^{(\zeta-1)}} - \theta_{q_2^{(\zeta+1)}}}{2} \right)}
   {\sin\left( \frac{\theta_{q_2^{(\zeta-1)}} - \theta_{q_2^{(\zeta)}}}{2} \right)\sin\left( \frac{\theta_{q_2^{(\zeta)}} - \theta_{q_2^{(\zeta+1)}}}{2} \right)\sin\left( \frac{\theta_{q_1^{(\xi-1)}} - \theta_{q_1^{(\xi+1)}}}{2} \right)}\cr
    &=\frac{\sin\left( \frac{\theta_{q_1^{(\xi-1)}} - \theta_{q_1^{(\xi)}}}{2} \right)\sin\left( \frac{\theta_{q_1^{(\xi)}} - \theta_{q_1^{(\xi+1)}}}{2} \right)\sin\left( \frac{\theta_{q_2^{(\zeta-1)}} - \theta_{q_2^{(\zeta+1)}}}{2} \right)}
   {\sin\left( \frac{\theta_{q_2^{(\zeta-1)}} - \theta_{q_2^{(\zeta)}}}{2} \right)\sin\left( \frac{\theta_{q_2^{(\zeta)}} - \theta_{q_2^{(\zeta+1)}}}{2} \right)\sin\left( \frac{\theta_{q_1^{(\xi-1)}} - \theta_{q_1^{(\xi+1)}}}{2} \right)}
   \cdot \frac{\sin\left( \frac{\theta_{q_1^{(\xi-1)}} - \theta_{q_2^{(\zeta-1)}}}{2} \right)}{\sin\left( \frac{\theta_{q_1^{(\xi-1)}} - \theta_{q_2^{(\zeta-1)}}}{2} \right)}\cr
    &=\frac{\sin\left( \frac{\theta_{N} - \theta_{q_1^{(\xi+1)}}}{2} \right)\sin\left( \frac{\theta_{q_1^{(\xi-1)}} - \theta_{q_2^{(\zeta-1)}}}{2} \right)}{\sin\left( \frac{\theta_{q_1^{(\xi-1)}} - \theta_{q_1^{(\xi+1)}}}{2} \right)\sin\left( \frac{\theta_{q_2^{(\zeta-1)}} - \theta_{N}}{2} \right)}\cdot
    \frac{\sin\left( \frac{\theta_{q_1^{(\xi-1)}} - \theta_{N}}{2} \right)\sin\left( \frac{\theta_{q_2^{(\zeta-1)}} - \theta_{q_2^{(\zeta+1)}}}{2} \right)}{\sin\left( \frac{\theta_{q_1^{(\xi-1)}} - \theta_{q_2^{(\zeta-1)}}}{2} \right)\sin\left( \frac{\theta_{N} - \theta_{q_2^{(\zeta+1)}}}{2} \right)}
    \cr
    &=-\langle N, q_1^{(\xi+1)}, q_1^{(\xi-1)}, q_2^{(\zeta-1)}\rangle\langle N, q_1^{(\xi-1)}, q_2^{(\zeta-1)}, q_2^{(\zeta+1)}\rangle.
\end{align}
Thus, we obtain
\begin{align}
    \Psi^N_{\bm{q}_1, \bm{q}_2}= -\Psi^{N-1}_{\bm{q}'_1, \bm{q}'_2}
    \langle N, q_1^{(\xi+1)}, q_1^{(\xi-1)}, q_2^{(\zeta-1)}\rangle\langle N, q_1^{(\xi-1)}, q_2^{(\zeta-1)}, q_2^{(\zeta+1)}\rangle.
\end{align}
\end{enumerate}
%
From the above results, which cover all possible cases,  $\Psi^N_{\bm{q}_1, \bm{q}_2}$ is reduced to $\Psi^{N-1}_{\bm{q}'_1, \bm{q}'_2}$ or decomposed as a product of $\Psi^{N-1}_{\bm{q}'_1, \bm{q}'_2}$ and invariants $\langle  i,j,k,l\rangle$.
\end{proof}

\end{document}
